\documentclass[12pt]{article}
\usepackage[utf8]{inputenc}
\usepackage[T1]{fontenc}
\usepackage{hyperref}
\title{Murut, Sembakung Language Resources for K-12 Curriculum Development}
\author{Compiled by Claude Sonnet 5.5 (claude-sonnet-5-5) for Andr\'as Kornai}
\date{}
\begin{document}
\maketitle

\section{Introduction}
Sembakung Murut (ISO 639-3: sbr; also Tinggalan) is an Austronesian language of the Murutic branch of North Borneo languages, spoken around the Sembakung river in Nunukan, North Kalimantan, Indonesia, and in Sabah, Malaysia. Estimates vary (about 3,200 speakers in Indonesia; Global Recordings lists 5,000). It is distinct from Selungai Murut and from other Murut languages such as Timugon, so materials for those must not be attributed to it. Searches found audio-visual Bible lessons and reference entries, but no dictionary, primer or school materials for Sembakung.

\section{Primary Educational Resources}
No Sembakung Murut primers, textbooks or school programmes were found. The language is not known to be taught in schools.

\section{Linguistic and Reference Materials}
\begin{itemize}
\item Wikipedia, Sembakung language: \url{https://en.wikipedia.org/wiki/Sembakung_language}
\item Numeral data for Sembakung Murut (Max Planck Institute, Channumerals): \url{https://lingweb.eva.mpg.de/channumerals/Murut-Sembakung.htm}
\item Wikitongues language page: \url{https://wikitongues.org/languages/sbr/}
\end{itemize}

\section{Religious and Community Resources}
\begin{itemize}
\item Joshua Project, Sembakung Murut in Indonesia (scripture status: portions): \url{https://www.joshuaproject.net/index.php/people_groups/14841/ID}
\item Joshua Project, Sembakung Murut in Malaysia: \url{https://www.joshuaproject.net/index.php/people_groups/14841/MY}
\end{itemize}

\section{Digital and Audio Resources}
Global Recordings Network offers ``Aho Onsoi'' (Good News) audio-visual Bible lessons, with an MP3 audio download and slideshow: \url{https://globalrecordings.net/en/language/sbr}

\section{Curriculum Development Context}
Materials for a Sembakung Murut curriculum essentially do not yet exist in the sources found. Development would need community-led orthography, dictionary and reading material work. A Webonary dictionary for the related Timugon Murut was reported (\url{https://www.malaymail.com/amp/news/life/2026/01/02/a-language-carried-forward-timugon-murut-dictionary-enters-the-digital-age/204040}). It is a different language, but it could serve as a model for tooling. Indonesian and Malay are the languages of schooling.
\end{document}